\chapter{Тестирование и апробация системы} \label{ch3}

В рамках данной главы проводится проверка работоспособности разработанной системы мониторинга, а также её апробация в условиях, приближенных к реальной эксплуатации. Основной целью тестирования является оценка корректности сбора и визуализации метрик, а также проверка функционирования механизма оповещений при возникновении различных типов нагрузок и нештатных ситуаций.

\section{Описание тестовой среды} \label{ch3:sec1}

В качестве тестовой базы данных была использована синтетическая база данных стандарта TPC-H \cite{tcph}, предназначенная для моделирования аналитической нагрузки. Данный набор данных широко применяется для оценки производительности систем управления базами данных и позволяет воспроизводить сложные сценарии обработки запросов.

На основе TPC-H была развернута база данных в СУБД Greenplum, содержащая набор таблиц, отражающих бизнес-процессы (заказы, поставщики, клиенты, регионы и др.). Для проведения тестирования использовался стандартный набор из 22 аналитических SQL-запросов (Q1–Q22), каждый из которых характеризуется значительным объёмом вычислений, множественными соединениями таблиц и операциями агрегации.

\section{Нагрузочное тестирование системы}  \label{ch3:sec2}

Нагрузочное тестирование проводилось путём последовательного выполнения набора TPC-H запросов. Для этого был сформирован единый SQL-скрипт, содержащий все 22 запроса, выполнение которого осуществлялось с использованием утилиты \texttt{psql}. При запуске скрипта дополнительно фиксировалось время выполнения каждого запроса, что позволило оценить их ресурсоёмкость.

В ходе тестирования рассматривался сценарий последовательного выполнения запросов, позволяющий оценить базовую производительность системы и характер нагрузки, создаваемой отдельными аналитическими операциями. Такой подход обеспечивает контролируемые условия эксперимента и позволяет проследить влияние каждого запроса на состояние системы мониторинга и значения метрик.

В процессе выполнения запросов осуществлялось наблюдение за показателями, отображаемыми на дашборде Grafana. На \firef{fig:test_1}, \firef{fig:test_2} и \firef{fig:test_3} было зафиксировано, что при увеличении интенсивности нагрузки наблюдается:

\begin{itemize}
	\item значительный рост загрузки процессора (на \firef{fig:test_3});
	\item увеличение потребления оперативной памяти;
	\item повышение интенсивности операций ввода-вывода;
	\item рост показателей транзакционной активности (TPS) (на \firef{fig:test_2}).
\end{itemize}

\begin{figure}[ht!] 
	\center
	\includegraphics [scale=0.15] {/Users/vadimvasilev/Study/8 семестр/НИР/latex/SPbPU-student-thesis-template/my_folder/images//test_1}
	\caption{Показатели блока "Greenplum Overview" при выполнении тестовой нагрузки} 
	\label{fig:test_1}  
\end{figure}

\begin{figure}[ht!] 
	\center
	\includegraphics [scale=0.15] {/Users/vadimvasilev/Study/8 семестр/НИР/latex/SPbPU-student-thesis-template/my_folder/images//test_2}
	\caption{Показатели блока "Database Metrics" при выполнении тестовой нагрузки} 
	\label{fig:test_2}  
\end{figure}

\begin{figure}[ht!] 
	\center
	\includegraphics [scale=0.15] {/Users/vadimvasilev/Study/8 семестр/НИР/latex/SPbPU-student-thesis-template/my_folder/images//test_3}
	\caption{Показатели блока "Database Metrics" при выполнении тестовой нагрузки} 
	\label{fig:test_3}  
\end{figure}

Наиболее ресурсоёмкие запросы (например, содержащие множественные соединения и агрегации) приводили к заметным пиковым значениям загрузки CPU и дисковой подсистемы. Это подтверждает корректность работы системы мониторинга, так как изменения метрик адекватно отражают характер выполняемой нагрузки.

%\begin{figure}[ht!]

%\end{figure}

\section{Проверка системы оповещений}  \label{ch3:sec3}

Для апробации механизма оповещений были проведены тесты, направленные на искусственное создание ситуаций, при которых должны срабатывать заданные правила.

В ходе тестирования были рассмотрены следующие сценарии:

\begin{itemize}
	\item Высокая загрузка процессора — создавалась за счёт выполнения ресурсоёмких запросов. При превышении порогового значения загрузки CPU срабатывало оповещение \textit{HighCPUUsage};
	
	\item Длительно выполняющиеся запросы — достигались за счёт выполнения сложных аналитических запросов. Это приводило к срабатыванию оповещения \textit{LongRunningQueries};
	
	\item Недоступность компонентов — моделировалась остановкой одного из контейнеров (например, экспортера), что приводило к срабатыванию оповещения \textit{ExporterDown}.
\end{itemize}

При наступлении указанных условий система Prometheus формировала события оповещений, которые передавались в Alertmanager. Далее происходила их обработка, группировка и отправка уведомлений на заданный адрес электронной почты.

%\begin{figure}[ht!]

%\end{figure}

В ходе тестирования на \firef{fig:alert_1}, \firef{fig:alert_2} и \firef{fig:alert_3} было подтверждено, что уведомления формируются корректно, содержат необходимую информацию о типе события, уровне критичности и описании проблемы, а также доставляются пользователю с минимальной задержкой.

\begin{figure}[ht!] 
	\center
	\includegraphics [scale=0.15] {/Users/vadimvasilev/Study/8 семестр/НИР/latex/SPbPU-student-thesis-template/my_folder/images//alert_1}
	\caption{Пример оповещения \textit{HighCPUUsage} полученный на электронную почту} 
	\label{fig:alert_1}  
\end{figure}

\begin{figure}[ht!] 
	\center
	\includegraphics [scale=0.15] {/Users/vadimvasilev/Study/8 семестр/НИР/latex/SPbPU-student-thesis-template/my_folder/images//alert_2}
	\caption{Пример оповещения \textit{ExporterDown} полученный на электронную почту} 
	\label{fig:alert_2}  
\end{figure}


\begin{figure}[ht!] 
	\center
	\includegraphics [scale=0.15] {/Users/vadimvasilev/Study/8 семестр/НИР/latex/SPbPU-student-thesis-template/my_folder/images//alert_3}
	\caption{Пример оповещения \textit{LongRunningQueries} полученный на электронную почту} 
	\label{fig:alert_3}  
\end{figure}
\newpage 
\newpage
\section{Анализ результатов тестирования}  \label{ch3:sec4}

Проведённые испытания показали, что разработанная система мониторинга корректно выполняет все основные функции:

\begin{itemize}
	\item обеспечивает стабильный сбор метрик в условиях различной нагрузки;
	\item корректно отражает изменения состояния системы на дашборде;
	\item позволяет выявлять узкие места и ресурсоёмкие операции;
	\item своевременно реагирует на критические события посредством системы оповещений.
\end{itemize}

\section{Выводы по главе} \label{ch3:conclusion}

Использование TPC-H запросов в качестве нагрузки позволило продемонстрировать поведение системы в условиях, приближенных к реальным аналитическим сценариям. Это подтверждает применимость разработанного решения для мониторинга производительности СУБД Greenplum в практических задачах.

Таким образом, результаты тестирования свидетельствуют о работоспособности, устойчивости и практической значимости разработанной системы мониторинга.
	




%% Вспомогательные команды - Additional commands
%
%\newpage % принудительное начало с новой страницы, использовать только в конце раздела
%\clearpage % осуществляется пакетом <<placeins>> в пределах секций
%\newpage\leavevmode\thispagestyle{empty}\newpage % 100 % начало новой страницы